\documentclass[11pt,a4paper]{article}
\usepackage[utf8]{inputenc}
\usepackage[spanish]{babel}
\usepackage[margin=1in]{geometry}
\usepackage{hyperref}
\usepackage{enumitem}
\usepackage{xcolor}

% Configuración de colores (Estilo sobrio e industrial)
\definecolor{primary}{RGB}{17, 17, 17}
\definecolor{accent}{RGB}{230, 80, 0}
\definecolor{muted}{RGB}{100, 100, 100}

\hypersetup{
    colorlinks=true,
    linkcolor=accent,
    urlcolor=accent,
}

\begin{document}
\pagestyle{empty}

% Cabecera
\begin{center}
    {\Huge \textbf{\textsc{Adrián Brihuega Sánchez}}} \\ \vspace{4pt}
    {\color{accent} \large \textbf{Ingeniería de Sistemas y Computación Estructurada}} \\ \vspace{6pt}
    {\color{muted} Madrid, España \quad $\cdot$ \quad \href{mailto:adrianbrihuegasanchez20@gmail.com}{adrianbrihuegasanchez20@gmail.com} \quad $\cdot$ \quad \href{https://github.com/brihuaa}{github.com/brihuaa}}
\end{center}

\vspace{15pt}

% Perfil
\section*{Perfil Técnico}
\hrule
\vspace{6pt}
Estudiante de último curso de Ingeniería Informática con sólidas bases en arquitectura de computadores, diseño lógico de hardware (VHDL), concurrencia y desarrollo nativo utilizando C/C++. Especializado en la optimización de flujos de sistemas, análisis estadístico clásico mediante sintaxis avanzada en SPSS y gestión de infraestructura de red en entornos Unix/Linux sin dependencias de cajas negras.

\vspace{10pt}

% Educación
\section*{Educación}
\hrule
\vspace{6pt}
\textbf{Grado en Ingeniería Informática} \hfill 2022 -- 2026 \\
Universidad Antonio de Nebrija \\
\textit{Enfoque principal:} Arquitectura de Computadores, Sistemas Operativos, Estructuras de Datos, Redes y Optimización de Código Basada en Hardware.

\vspace{10pt}

% Habilidades
\section*{Entorno Técnico}
\hrule
\vspace{6pt}
\begin{itemize}[leftmargin=*, noitemsep]
    \item \textbf{Bajo Nivel y Tipados:} C++, C, Java, VHDL, Arquitectura de Hardware (Xilinx).
    \item \textbf{Sistemas y Lógica:} Unix Kernel, Linux Bash, Administración de Redes, SPSS (Sintaxis Avanzada).
    \item \textbf{Infraestructura y Backend:} Node.js, TypeScript, REST APIs, SQL, Docker CLI, Git Server.
\end{itemize}

\vspace{10pt}

% Proyectos
\section*{Proyectos Destacados}
\hrule
\vspace{6pt}

\textbf{Core-Scheduler C++} \\
\textit{Planificador de tareas multinivel simulado en C++ de alto rendimiento. Optimiza colas de procesos utilizando algoritmos Round Robin adaptativos, monitorizando el impacto y las fallas en la caché de la CPU mediante Valgrind.}

\vspace{6pt}

\textbf{VHDL Processor-Unit} \\
\textit{Diseño e implementación de una ALU monociclo de 8 bits completa utilizando VHDL sobre suites de desarrollo Xilinx. Testeada e integrada a nivel físico estructural.}

\vspace{6pt}

\textbf{Statistical Labor Cost Pipeline} \\
\textit{Procesamiento e ingesta masiva de series temporales de datos del INE (2008--2024) utilizando sintaxis automatizada en SPSS para depuración y análisis de regresión lineal.}

\end{document}